\scanpage{232}
\mspage{106}
% archive sheet number: 231

de $A^\wedge$ (ou, ce qui revient au $=$, \deleted{que} de $(A_{/a})^\wedge$).

\inserted{$c)$ Il existe un objet $I$ de $A^\wedge$ tel que $I\to e_A$ soit
$W_A$-asphérique, \deleted{et que} avec deux sections $\delta_0,\delta_1$ de $I$ sur
$e_A$, \unsure{avec|telles} $\mathrm{Ker}(\delta_0,\delta_1)=\varnothing_A$ (préfaisceau vide).}

$d)$ Les catégories localisées $A_{/a}=B$ sont des
\uline{catégories-test} pour $W$ (cf Pursuing Stacks, \S\,65, notamment pages 172--173)
ce qui revient à dire que l'on a ceci pour chacune des \deleted{\unsure{ca}} catégories $B=A_{/a}$ :

% slanted
\margin{\uline{NB} La condition $d)$, d'apparence très forte, est équivalente à la condition d'\unsure{apparence|allure} plus \unsure{anodine} $d')$ Le foncteur $j_A:\mathrm{Cat}\to A^\wedge$ transforme \unsure{flèches|éléments} à \unsure{fct adjoint|fonc. adjoints} en flèches $W_A$-asphériques. -- et il suffit $=$ de l'exiger pour $C\to e$ dans Cat, où $C$ \deleted{est} un objet de Cat ayant objet final.}

\deleted{(8.4.2)} Les foncteurs

\[
 \left\{\begin{aligned}
 i_B &: B^\wedge\longrightarrow\mathrm{Cat}
 & (F\unsure{\mapstochar\relbar|\longmapsto} B_{/F})\\
 j_B &: \mathrm{Cat}\unsure{\mathrel{\text{---}}|\longrightarrow} B^\wedge
 &\deleted{(C\longmapsto\mathrm{Hom}_{\mathrm{Cat}}(}\\
 &&(C\mapsto(b\longmapsto\mathrm{Hom}(B_{/b},C)))
 \end{aligned}\right.
 \tag*{(8.4.2)}
\]

\deleted{qui} où $j_B$ est l'adjoint à droite de $i_B$) sont tous les deux ``model preserving'' i.e.

\[
 \left\{\begin{aligned}
 W_B&=i_B^{-1}(W) &&\text{(par définition de }W_B\text{ en fait)}\\
 W&=j_B^{-1}(W_B)
 \end{aligned}\right.
 \tag*{(8.4.3.)}
\]

et les foncteurs \deleted{local} localisés

\diagram{\begin{tikzcd}[ampersand replacement=\&]
 \mathrm{Hot}_{B,W}\arrow[r,shift left,"\overline{i_B}"] \&
 \mathrm{Hot}_{W}\arrow[l,shift left,"\overline{j_B}"]
\end{tikzcd}}

sont des équivalences de catégories. De plus, ces équivalences de catégories sont
quasi-inverses l'une de l'autre, les morphismes d'adjonction étant déduits
des morphismes d'adjonction pour les couples foncteurs \inserted{adj.} $(i_B,j_B)$, ces morphismes d'adjonction
étant des \unsure{flèches|élts} dans $W$ respt. dans $W_A$.
